% !TEX root = ../../../main.tex
% 来源：中学数学实验教材·第二册（几何）
% 章节：直线形 / 平行线与内角和定理
% 原始文件：3.tex，行 2959-2987
% 类型：tikz
% ---
\begin{tikzpicture}
\begin{scope}
    \tkzDefPoints{-2/0/B, 2/0/C, 0/2/A, .5/0/D, .5/.5/G}
\tkzDrawPolygon(A,B,C)
\tkzDefPointWith[linear, K=1.5](B,A) \tkzGetPoint{H}
\tkzDrawSegments[dashed](A,H)
\tkzInterLL(B,A)(D,G) \tkzGetPoint{E}
\tkzInterLL(A,C)(D,E) \tkzGetPoint{F}
\tkzLabelPoints[below](B,C,D)
\tkzLabelPoints[above](A,E)
\tkzLabelPoints[right](F)
\tkzDrawSegments(D,E)
\end{scope}
\begin{scope}[xshift=5cm, scale=.6]
    \tkzDefPoints{-2/0/B, 2/0/C, 0/2/A, 3/0/D, 3/.5/G}
    \tkzDrawPolygon(A,B,C)
    \tkzDefPointWith[linear, K=3](B,A)  \tkzGetPoint{H}
    \tkzDefPointWith[linear, K=1.5](B,C)  \tkzGetPoint{I}
    \tkzDefPointWith[linear, K=1.5](A,C)  \tkzGetPoint{J}
    \tkzInterLL(B,A)(D,G) \tkzGetPoint{E}
    \tkzInterLL(A,C)(D,E) \tkzGetPoint{F}
    \tkzLabelPoints[below](B,C)
    \tkzLabelPoints[above](A,E)
    \tkzLabelPoints[right](F)
    \tkzLabelPoints[above right](D)
    \tkzDrawSegments[dashed](A,H C,J C,I)
    \tkzDrawSegments(E,F)
\end{scope}
\end{tikzpicture}
